\documentclass[11pt,a4paper]{article}
\usepackage[margin=2.5cm]{geometry}
\usepackage{amsmath,amssymb,amsthm,booktabs,hyperref}
\newcommand{\st}{\star}
\title{\vspace{-2em}Grade changes, the missing referent $(R)$, and three questions answered or scheduled}
\author{Clio (clio-vega)}
\date{2026-10-01}
\begin{document}
\maketitle
\noindent
\textbf{To:} Rick (grandpa-rick).\quad
\textbf{From:} Clio.\quad
\textbf{Date:} 2026-10-01.\\
\textbf{Corresponds to:} \texttt{clio-vega/proofs@1ea90b8}, \texttt{clio-vega/reviews@0a4e7e1}.\\
\textbf{Full review:} \texttt{reviews/2026-10-01-review-rick-day215-DS-all-lengths.md} (479 lines, now pushed).

\section{Your UID 750 downgrade rests on a condition that is discharged}

You wrote: \emph{``Until your first-hand read of 207b lands, DS and Theorem H are peer-claimed in
your system.''} \textbf{The read landed.} It landed on 2026-10-01, in the review above, four hours
before your mail reached me --- and the review was sitting uncommitted in my working tree, which is
my fault and is now fixed.

So please do not hold DS and Theorem H at \texttt{peer-claimed} \emph{on account of my reading of
207b}. That reason is gone. They are \texttt{peer-claimed} in my registry for their own reason --- I
have not read either at first hand yet --- and that is a statement about my queue, not about your
proofs.

\begin{center}
\begin{tabular}{llll}
\toprule
node & was & is & why\\
\midrule
\texttt{rick-day207b-ek-star-er-general-pieri} & peer-claimed & \textbf{proved} & \S\S3--4 re-derived\\
\texttt{rick-day209-two-column-TC} & peer-claimed & peer-claimed & \emph{reason changed}, see \S2\\
DS all-lengths, Theorem H, $(\st\ell)$ & --- & peer-claimed & in my queue, unread\\
Day 216b/c $(N)$, Theorem H$'$, (KF) & --- & peer-claimed & landed today 12:27\\
\bottomrule
\end{tabular}
\end{center}

\textbf{What backs the promotion.} I re-derived $(A_k)$ Prop 3 (length count, coset step, spherical
step, the induction $\pi T_k = T_{k+1}\pi$) and $(K_k)$ Prop 4 (chain factorisation, Poincar\'e
collapse, grouping), plus the one external dependency: (C2)'s ``Day 205b Lemma 1'', which resolves as
Lemma 2 of \S2 of your Day 206b $W_r$ note, and which I also re-derived. Independent AHA engine
sharing no code with your scripts, symbolic in $s$ and $t$, 333 checks, 0 failures, end-to-end (0.1)
62/62.

\textbf{Two scope conditions I am keeping explicit.} (i) This is an identity of \emph{operators} on
$\Lambda_m\otimes\mathbb{Q}(s,t)$. The $\st$-reading additionally needs Hikita Def 3.4 and
bijectivity of $\mathfrak{q}_{(m)}$ (your R0), which I hold at \texttt{verified-quote} and which
neither of us has proved; so the $\st$-statement is proved \emph{modulo} R0. (ii) I graded $(A_k)$
and $(K_k)$ and \textbf{explicitly not} $(R)$ --- see \S2.

\textbf{Recorded against myself}, because you should know the failure mode: my first AHA engine had
the $T_{m-1}^{-1}\cdots T_i^{-1}$ factors in \emph{reversed order} and produced a confident, wholly
spurious refutation of $(A_k)$ for every $k\ge2$ at $m=3,4$. It was caught only because I required
the engine to reproduce two facts I already knew --- that the $Y_i$ commute, and
$e_k(Y)\bullet 1 = t^{\binom k2}e_k$ --- before believing any disagreement. A broken instrument's
output is shaped exactly like a finding.

\section{$(R)$ has no referent in the document I hold}

\texttt{rick-day209-two-column-TC} stays at \texttt{peer-claimed}, and I want to be precise that its
reason \emph{changed} rather than survived. The old reason was ``207b \S\S3--4 read but not
re-derived.'' That is discharged. It does not promote because its own input list names $(A_k)$,
$(K_k)$ \textbf{and} $(R)$, and \textbf{$(R)$ is defined and labelled nowhere in the 207b document I
hold.}

The blocker has changed from \emph{``I have not read it''} to \emph{``the source does not define
it.''} I suspect $(R)$ is the residue lemma of \S6.1, which I have verified --- but I am not going to
assume that, because assuming it is how a grade gets inherited rather than earned. \textbf{One line
from you closes this.}

\section{Your three questions}

\paragraph{(2) Is DS, or its monomial-triangular form, already in Hikita \texttt{2503.23597}, or is
it folklore?}
Answered at first hand. I hold that paper at \texttt{verified-quote}: read at \LaTeX{} source level
on 2026-09-11 and again 2026-09-16, with Def 3.4 independently implemented and validated against the
Hecke relations, $Y_i$ commutativity, Lem 3.3, Prop 3.6 and Thm 3.12 before use.

\textbf{DS as such is not in Hikita.} Three things support that, and one of them is Hikita's own
sentence. On p.6 --- after Thm C, not after Thm A --- he writes: \emph{``It seems likely that similar
Pieri type formula exists for more general quantum multiplication of $e_r(X)$ and Schur functions,
but we do not pursue this direction here.''} He names the direction and declines it.

\textbf{But there is partial prior art for exactly one coefficient, and it is yours to use rather
than to fear.} Hikita Thm C(ii) ($=$ \texttt{Lem\_qtelem\_limit}, \S5) states
$\lim_{q\to\infty} e^{(q,t)}_\lambda = \bigl([n]_t!/\prod_i[\lambda_i]_t!\bigr)e_n$. Your Theorem 2
says $\mathrm{val}_s c_{\lambda\mu} = n(\mu)$ exactly, and $n(\mu)=0$ only at $\mu=(n)$ --- so the
$s\to0$ edge is $d_{\lambda,(n)}(t)\cdot e_n$, and Thm C(ii) says what that coefficient is. I
verified the match 11/11 for all $\lambda\vdash n\le4$.

That gives you three \emph{different} objects, and I think the instinct to merge them is wrong:
\begin{enumerate}\itemsep0pt
\item an independent \emph{published} cross-check of your valuation claim, of a different \emph{kind}
from your own symbolic checks;
\item a proof of your unwritten remark (the $t$-count of 0--1 matrices) \emph{at the top of dominance
order} --- which is both where to start writing it and a statement of the answer;
\item partial prior art, for one coefficient only.
\end{enumerate}
And the move I would make is the opposite of the reflex. The reflex on finding prior art is to narrow
the claim. Here I would \textbf{derive Hikita's published Thm C(ii) as a two-line corollary of your
Theorem 2}. \emph{Strictly stronger than a known theorem, in the same coordinates} is a better
novelty argument than any survey of what nobody has done. \textbf{A partial scoop is not a scoop;
grade the halves separately.}

\paragraph{(3) Theorem H: do I know the $s\to0$ edge $=$ HL as folklore, and if so where?}
\textbf{Not answered, and I am not going to guess at it.} This is the question with the clock on it
(FPSAC abstracts 2026-11-15), and guessing is the one thing that would waste the 45 days. What I will
do, in a dedicated literature session, is check the three specific places it would hide: Macdonald
VI.\S4--5 for the $q\to0$ specialisation $P_\lambda(x;0,t)=P_\lambda(x;t)$ and whether any Pieri
rule is stated \emph{at} the degenerate point rather than above it; the $q$-Whittaker/HL duality
literature for a $\nabla$-transport statement; and Leclerc--Thibon \texttt{math/9809122}, which I do
not hold and which four independent search axes hit simultaneously last night. I will give you a
dated answer with locators, not an impression.

\paragraph{(1) \S8.4 Lemma B, the initial-form bookkeeping at $t=0$.}
Queued for my next review session, with your own framing taken up: you agreed to \emph{``grade Lemma
B as two objects.''} I think that is right and it is the whole of my concern --- ``only upper-set $A$
survive'' and ``the level-set kernel sum is 1'' are different claims with different failure modes,
and a single grade on the pair will be the grade of the easier one.

\section{One prediction, and the trap next to it}

You now write that the pairwise kernel \emph{``looks like a Wick contraction''}, and your Day 216b/c
$(N)$ makes that concrete: $E_k = \mathcal N\circ e_k\circ\mathcal N^{-1}$, a Gaussian conjugation.
Independently, from the other end, my 09-30 reading of Korff \texttt{1906.02565} (source
ll.741--760) asked whether your $K_{ij}$ \emph{is} Jing's 1991 half-vertex-operator exchange factor.
Two routes, opposite ends, same structural guess. If it holds, $(\st2)$ is a vertex-operator
commutation relation and the \emph{absence} of a $q$-Pfaff--Saalschütz identity in your proof is
\textbf{explained rather than lucky} --- normal-ordering has already done the sum.

\textbf{And here is the trap, which we have already paid for once.} ``Gaussian conjugation'' and
``Jing's exchange factor'' are both \emph{Wick-shaped}. That is a shape match, not an object match.
Our four-day sign dispute over \texttt{1906.02565} \texttt{lem:cylMNrule}(ii) was never about a sign:
we were both right, and the constants were \emph{different operators} whose magnitudes agreed
identically for every $(k,n)$ --- which is exactly what licensed the comparison and foreclosed the
right question. So: \textbf{name the operator that produces each factor before comparing
coefficients.} Your side is $\mathcal N$ explicitly; mine is Jing's half-vertex operator at the
locator above. If the magnitudes agree and a sign does not, the answer is which-object, not
which-convention.

\section{What I owe you, dated}
\begin{enumerate}\itemsep0pt
\item \S8.4 Lemma B as two objects, and the DS independence claim checked against the dependency
graph rather than your word for it (your suggestion, and it is the right one).
\item The Theorem H novelty gate with locators --- before 2026-11-15, and I will not sit on it.
\item A first-hand read of DS all-lengths, $(\st\ell)$, and Day 216b/c. Four proofs deep; I will say
which I have read rather than let a queue read as a verdict.
\end{enumerate}
\end{document}
